\documentclass[../../../include/open-logic-section]{subfiles}

\begin{document}

\olfileid[id]{sth}{choice}{countablechoice}
\olsection{Pilihan Terhitung}

Mudah dibuktikan, tanpa menggunakan Pilihan/Pengurutan Baik sama sekali,
bahwa:

\begin{lem}[dalam $\Zminus$]
Setiap himpunan hingga mempunyai suatu fungsi pilihan.
\end{lem}

\begin{proof}
Misalkan $a = \{b_1, \ldots, b_n\}$. Demi kesederhanaan, andaikan setiap
$b_i\neq\emptyset$. Jadi, terdapat objek-objek $c_1,\ldots,c_n$ sedemikian
sehingga $c_1\in b_1,\ldots,c_n\in b_n$. Berdasarkan
\olref[z][pairs]{prop:pairsconsequences}, himpunan
$\{\langle b_1,c_1\rangle,\ldots,\langle b_n,c_n\rangle\}$ ada; himpunan
ini merupakan suatu fungsi pilihan bagi~$a$.
\end{proof}

Namun, keadaan menjadi lebih keruh begitu kita mempertimbangkan
himpunan-himpunan tak hingga. Sebagai contoh, pertimbangkan perluasan
``minimal'' berikut dari hasil di atas:

\begin{defish}
\emph{Pilihan Terhitung.} Setiap himpunan \emph{terhitung} mempunyai suatu
fungsi pilihan.
\end{defish}

Ini merupakan suatu kasus khusus Pilihan. Ternyata prinsip ini cukup sering
digunakan tanpa kesadaran yang jelas akan penggunaannya. Berikut dua contoh
yang bagus.\footnote{Berasal dari \citet[\S9.4]{Potter2004} dan Luca
Incurvati.}

\begin{ex}
Berikut suatu gagasan alami: bagi sebarang himpunan $A$, berlaku
$\cardle{\omega}{A}$ atau $\cardeq{A}{n}$ bagi suatu $n\in\omega$. Ini
merupakan salah satu cara untuk menyatakan gagasan intuitif bahwa setiap
himpunan bersifat hingga atau tak hingga. Cantor dan banyak matematikawan
lain mengemukakan klaim ini tanpa membuktikannya. Dengan kehati-hatian kita,
kita telah membuktikannya dalam
\olref[cardinals][classing]{generalinfinitycharacter}. Namun, dalam bukti itu
kita bekerja dalam $\ZFC$ karena mengandaikan bahwa setiap himpunan $A$
dapat diurutkan baik, sehingga keberadaan $\card{A}$ terjamin. Artinya, kita
secara eksplisit mengandaikan Pilihan.

Sesungguhnya, \citet{Dedekind1888} memberikan buktinya sendiri bagi klaim
ini sebagai berikut:

\begin{thm}[dalam $\Zminus + \text{Pilihan Terhitung}$]
Untuk sebarang $A$, berlaku $\cardle{\omega}{A}$ atau $\cardeq{A}{n}$ bagi
suatu $n \in \omega$.
\end{thm}

\begin{proof}
Andaikan $\cardneq{A}{n}$ bagi semua $n\in\omega$. Maka, khususnya, bagi
setiap $n<\omega$ terdapat himpunan bagian $A_n\subseteq A$ dengan tepat
$\cardexpo{2}{n}$ anggota. Dengan menggunakan barisan
$A_0,A_1,A_2,\ldots$, bagi setiap $n$ kita definisikan:
\[
	B_n = A_n \setminus \bigcup_{i < n} A_i.
\]
Sekarang perhatikan:
\begin{align*}
	\card{\bigcup_{i < n}A_i}
	&\leq \card{A_0} + \card{A_1} + \ldots + \card{A_{n-1}}\\
	&=1 + 2 + \ldots + 2^{n-1}\\
	& = 2^n - 1\\
	& < 2^n = \card{A_n}
\end{align*}
Karena itu, setiap $B_n$ mempunyai sekurang-kurangnya satu anggota,
$c_n$. Selain itu, himpunan-himpunan $B_n$ saling lepas; jadi, jika
$c_n=c_m$, maka $n=m$. Namun, setiap $c_n\in A$. Dengan demikian, fungsi
$f(n)=c_n$ merupakan suatu injeksi $\omega\to A$.
\end{proof}
\noindent
Dedekind tidak menandai bahwa ia telah menggunakan Pilihan Terhitung. Namun,
\emph{apakah Anda} melihat penggunaannya? Perhatikan kembali.
(Sungguh: \emph{perhatikan kembali}.)

Bukti tersebut menggunakan Pilihan Terhitung dua kali. Kita menggunakannya
sekali untuk memperoleh barisan himpunan $A_0,A_1,A_2$, \dots\@ lalu sekali
lagi untuk memilih anggota $c_n$ dari setiap~$B_n$. Selain itu, penggunaan
Pilihan ini tidak dapat dihilangkan. \citet[hlm.~138]{Cohen1966}
membuktikan bahwa hasil tersebut gagal jika kita sama sekali tidak mempunyai
suatu bentuk Pilihan. Artinya, konsisten dengan $\ZF$ bahwa terdapat
himpunan-himpunan yang \emph{tidak dapat dibandingkan} dengan~$\omega$.
\end{ex}

\begin{ex}
Pada \citeyear{Cantor1878}, Cantor menyatakan bahwa suatu gabungan terhitung
dari himpunan-himpunan terhitung bersifat terhitung. Ia tidak menyajikan
bukti, barangkali karena menganggap bukti tersebut jelas. Dengan
kehati-hatian kita, kita telah membuktikan versi yang lebih umum dari hasil
ini dalam \olref[card-arithmetic][simp]{kappaunionkappasize}. Namun, bukti
kita secara eksplisit mengandaikan Pilihan. Bahkan bukti hasil yang kurang
umum pun memerlukan Pilihan Terhitung.

\begin{thm}[dalam $\Zminus + \text{Pilihan Terhitung}$]
Jika $A_n$ terhitung bagi setiap $n\in\omega$, maka
$\bigcup_{n<\omega}A_n$ terhitung.
\end{thm}

\begin{proof}
Tanpa mengurangi keumuman, andaikan setiap $A_n\neq\emptyset$. Jadi, bagi
setiap $n\in\omega$, terdapat !!a{surjection}
$f_n\colon\omega\to A_n$. Definisikan
$f\colon\omega\times\omega\to\bigcup_{n<\omega}A_n$ dengan
$f(m,n)=f_n(m)$. Hasilnya mengikuti karena $\omega\times\omega$ terhitung
(\olref[sfr][siz][zigzag]{natsquaredenumerable}) dan $f$ merupakan
!!a{surjection}.
\end{proof}
\noindent
Apakah Anda melihat penggunaan Pilihan Terhitung? Prinsip itu digunakan
untuk memilih barisan fungsi $f_0,f_1,f_2,\dots$.\footnote{Penggunaan Pilihan
yang serupa terjadi dalam
\olref[card-arithmetic][simp]{kappaunionkappasize}, ketika kita memberikan
instruksi ``Untuk setiap $\beta\in\cardfont{a}$, tetapkan
!!a{injection} $f_\beta$''.} Sekali lagi, hasil tersebut gagal jika tidak
ada prinsip Pilihan sama sekali. Secara khusus,
\citet{FefermanLevy1963} membuktikan bahwa konsisten dengan $\ZF$ bahwa
suatu gabungan terhitung dari himpunan-himpunan terhitung mempunyai
kardinalitas~$\beth_1$. Berikut pernyataan yang jauh lebih lucu mengenai
pokok ini dari Russell:
\begin{quote}
  Hal ini digambarkan oleh seorang jutawan yang membeli sepasang kaus kaki
  setiap kali ia membeli sepasang sepatu bot, dan tidak pernah pada waktu
  lain. Kegemarannya membeli keduanya begitu besar sehingga akhirnya ia
  mempunyai $\aleph_0$ pasang sepatu bot dan $\aleph_0$ pasang kaus
  kaki\dots\@ Di antara sepatu bot kita dapat membedakan yang kanan dan yang
  kiri, sehingga kita dapat memilih satu dari setiap pasang: kita dapat
  memilih semua sepatu bot kanan atau semua sepatu bot kiri. Namun, bagi kaus
  kaki tidak ada prinsip pemilihan seperti itu yang tersedia, dan tanpa
  mengandaikan aksioma multiplikatif [yakni, pada hakikatnya Pilihan], kita
  tidak dapat memastikan bahwa terdapat suatu kelas yang terdiri atas satu
  kaus kaki dari setiap pasang.
  \citep[hlm.~126]{Russell1919}
\end{quote}
Singkatnya, suatu bentuk Pilihan diperlukan untuk membuktikan hal berikut:
Jika Anda mempunyai pasangan kaus kaki dalam jumlah terhitung, maka Anda
mempunyai (hanya) sejumlah terhitung kaus kaki. Bahkan, tanpa Pilihan
Terhitung (atau sesuatu yang ekuivalen), gabungan terhitung dari
himpunan-himpunan terhitung dapat gagal bersifat terhitung.
\end{ex}

Pelajarannya ialah bahwa Pilihan Terhitung digunakan berulang kali tanpa
banyak disadari oleh para penggunanya. Pertanyaan filosofisnya ialah:
Bagaimana kita dapat \emph{membenarkan} Pilihan Terhitung?

Suatu upaya pembenaran intuitif mungkin mengacu pada suatu supertugas.
Andaikan kita membuat pilihan pertama dalam $\nicefrac{1}{2}$ menit,
pilihan kedua dalam $\nicefrac{1}{4}$ menit, \dots, pilihan ke-$n$ dalam
$\nicefrac{1}{2^n}$ menit, \dots\@ Dalam waktu $1$~menit, kita akan telah
membuat suatu barisan pilihan sepanjang $\omega$ dan mendefinisikan suatu
fungsi pilihan.

Namun, apa yang sebenarnya dapat ditunjukkan oleh eksperimen pikiran seperti
itu? Pertama, ia bergantung pada pemahaman yang sangat harfiah terhadap
gagasan ``memilih''. Selain itu, ia tampak mengikat matematika pada
kemungkinan metafisis.

Yang lebih penting, hal itu tidak akan memberikan pembenaran bagi Pilihan
\emph{secara penuh}, bukan \emph{sekadar} Pilihan Terhitung. Sebab, jika kita
memerlukan \emph{setiap} himpunan mempunyai suatu fungsi pilihan, kita perlu
mampu melakukan ``supertugas dengan panjang ordinal sebarang''. Terus
terang, gagasan itu menggelikan.

\end{document}
